\documentclass[aspectratio=169,11pt]{beamer}

\usepackage[T1]{fontenc}
\usepackage[utf8]{inputenc}
\usepackage{lmodern}
\usepackage{microtype}
\usepackage{booktabs}
\usepackage{tabularx}
\usepackage{array}
\usepackage{ragged2e}
\usepackage{xcolor}
\usepackage{hyperref}
\usepackage{amsmath}

\definecolor{StudentNavy}{HTML}{102A43}
\definecolor{StudentBlue}{HTML}{1769E0}
\definecolor{StudentCyan}{HTML}{14A6B8}
\definecolor{StudentGold}{HTML}{E8A628}
\definecolor{StudentInk}{HTML}{243B53}
\definecolor{StudentMist}{HTML}{F2F7FC}
\definecolor{StudentSoft}{HTML}{D9E8F5}
\definecolor{StudentGray}{HTML}{627D98}
\definecolor{StudentWhite}{HTML}{FFFFFF}
\definecolor{StudentGreen}{HTML}{2B8A6E}

\hypersetup{colorlinks=true,urlcolor=StudentBlue,linkcolor=StudentBlue}

\setbeamercolor{normal text}{fg=StudentInk,bg=StudentWhite}
\setbeamercolor{structure}{fg=StudentBlue}
\setbeamercolor{frametitle}{fg=StudentNavy,bg=StudentWhite}
\setbeamercolor{title}{fg=StudentWhite}
\setbeamercolor{subtitle}{fg=StudentSoft}
\setbeamercolor{author}{fg=StudentWhite}
\setbeamercolor{date}{fg=StudentSoft}
\setbeamercolor{block title}{fg=StudentWhite,bg=StudentBlue}
\setbeamercolor{block body}{fg=StudentInk,bg=StudentMist}

\setbeamerfont{title}{series=\bfseries,size=\fontsize{29}{34}\selectfont}
\setbeamerfont{subtitle}{size=\fontsize{14.5}{18}\selectfont}
\setbeamerfont{author}{series=\bfseries,size=\fontsize{15}{18}\selectfont}
\setbeamerfont{date}{size=\fontsize{11.5}{14}\selectfont}
\setbeamerfont{frametitle}{series=\bfseries,size=\fontsize{19.5}{23}\selectfont}
\setbeamerfont{normal text}{size=\fontsize{13.5}{17}\selectfont}
\setbeamerfont{block title}{series=\bfseries,size=\fontsize{13}{16}\selectfont}
\setbeamerfont{block body}{size=\fontsize{12}{15.5}\selectfont}
\setbeamerfont{footline}{size=\fontsize{8}{10}\selectfont}

\setbeamertemplate{navigation symbols}{}
\setbeamertemplate{items}[circle]
\setbeamertemplate{itemize item}{\color{StudentCyan}\raisebox{0.12ex}{\scriptsize$\blacktriangleright$}}
\setbeamertemplate{itemize subitem}{\color{StudentBlue}\scriptsize$\bullet$}
\setlength{\leftmargini}{1.2em}
\setlength{\leftmarginii}{1.1em}

\setbeamertemplate{frametitle}{%
  \nointerlineskip
  \begin{beamercolorbox}[wd=\paperwidth,leftskip=0.7cm,rightskip=0.7cm,sep=0.17cm]{frametitle}
    \rightskip=0pt plus 1fil\relax\usebeamerfont{frametitle}\insertframetitle\par
    \color{StudentCyan}\rule{\dimexpr\paperwidth-1.8cm\relax}{0.8pt}
  \end{beamercolorbox}
}

\setbeamertemplate{footline}{%
  \leavevmode
  \hbox{%
    \begin{beamercolorbox}[wd=.67\paperwidth,ht=2.7ex,dp=1.2ex,leftskip=.7cm]{author in head/foot}%
      \color{StudentGray}\usebeamerfont{footline}Sosialisasi Orang Tua \textbar\ SMAN 1 Muntilan
    \end{beamercolorbox}%
    \begin{beamercolorbox}[wd=.31\paperwidth,ht=2.7ex,dp=1.2ex,center]{date in head/foot}%
      \color{StudentBlue}\usebeamerfont{footline}Bab 2 \textbar\ \insertframenumber/\inserttotalframenumber
    \end{beamercolorbox}%
  }
}

\newcommand{\kicker}[1]{%
  {\color{StudentCyan}\bfseries\fontsize{10.5}{13}\selectfont\MakeUppercase{#1}}\par\vspace{0.11cm}
}

\newcommand{\claim}[1]{%
  {\color{StudentNavy}\bfseries\fontsize{18}{22}\selectfont #1}\par
}

\newcommand{\metric}[2]{%
  {\color{StudentBlue}\bfseries\fontsize{27}{30}\selectfont #1}\par
  \vspace{0.03cm}
  {\color{StudentGray}\fontsize{11.3}{14}\selectfont #2}
}

\newcommand{\answer}[1]{%
  \vspace{0.10cm}
  {\color{StudentGreen}\bfseries\fontsize{11.5}{14}\selectfont Jawaban: #1}
}

\newcommand{\statusok}{\textcolor{StudentGreen}{\bfseries RESMI 2026}}
\newcommand{\statuswait}{\textcolor{StudentGold}{\bfseries ANTISIPASI 2027}}

\title{Struktur, Materi, dan Contoh\\UTBK-SNBT}
\subtitle{Bab 2 \textbar\ Memahami tujuh subtes sebagai acuan pembimbingan 2027}
\author{Balya Rochmadi}
\date{Untuk Orang Tua/Wali Siswa \textbar\ Tahun Ajaran 2026/2027}

\begin{document}

% Slide 1 ---------------------------------------------------------------------
\begingroup
\setbeamercolor{background canvas}{bg=StudentNavy}
\begin{frame}[plain]
  \vspace{0.65cm}
  \kicker{SMAN 1 Muntilan \textbar\ StudentPro}
  \vspace{0.08cm}
  {\usebeamerfont{title}\color{StudentWhite}\inserttitle\par}
  \vspace{0.36cm}
  {\usebeamerfont{subtitle}\color{StudentSoft}\insertsubtitle\par}
  \vfill
  {\color{StudentGold}\rule{2.3cm}{2.5pt}}\par
  \vspace{0.20cm}
  {\usebeamerfont{author}\color{StudentWhite}\insertauthor}\par
  {\usebeamerfont{date}\color{StudentSoft}\insertdate}
  \vspace{0.42cm}
\end{frame}
\endgroup

% Slide 2 ---------------------------------------------------------------------
\begin{frame}{Bab ini membantu orang tua membaca kebutuhan belajar siswa}
  \kicker{Tiga Pertanyaan Utama}
  \begin{columns}[T,onlytextwidth]
    \begin{column}{0.31\textwidth}
      \begin{block}{Apa yang diujikan?}
        Dua kelompok tes, tujuh subtes, serta kemampuan yang berbeda pada setiap bagian.
      \end{block}
    \end{column}
    \begin{column}{0.31\textwidth}
      \begin{block}{Mengapa nilai berbeda?}
        Seorang siswa dapat kuat pada satu subtes dan memerlukan bantuan pada subtes lain.
      \end{block}
    \end{column}
    \begin{column}{0.31\textwidth}
      \begin{block}{Apa tindak lanjutnya?}
        Hasil per subtes digunakan untuk memilih materi, latihan, remedial, dan strategi waktu.
      \end{block}
    \end{column}
  \end{columns}
  \vspace{0.18cm}
  \centering
  {\color{StudentNavy}\bfseries\fontsize{13}{16}\selectfont Nilai total penting, tetapi profil kemampuan siswa jauh lebih berguna untuk pembimbingan.}
\end{frame}

% Slide 3 ---------------------------------------------------------------------
\begin{frame}{Framework resmi 2026 menjadi titik acuan sementara}
  \kicker{Struktur yang Sudah Terpublikasi}
  \begin{columns}[T,onlytextwidth]
    \begin{column}{0.20\textwidth}
      \metric{2}{kelompok tes}
    \end{column}
    \begin{column}{0.20\textwidth}
      \metric{7}{subtes}
    \end{column}
    \begin{column}{0.20\textwidth}
      \metric{155}{akumulasi soal}
    \end{column}
    \begin{column}{0.22\textwidth}
      \metric{195}{menit tanpa jeda}
    \end{column}
  \end{columns}
  \vspace{0.40cm}
  \begin{beamercolorbox}[rounded=true,sep=0.22cm,wd=\textwidth]{block body}
    \centering\bfseries Ketentuan UTBK 2027 belum diperlakukan final sebelum dipublikasikan oleh SNPMB.
  \end{beamercolorbox}
  \vspace{0.22cm}
  {\color{StudentGray}\fontsize{10.5}{13}\selectfont Angka 155 merupakan penjumlahan jumlah soal pada tujuh subtes dalam framework resmi 2026.}
\end{frame}

% Slide 4 ---------------------------------------------------------------------
\begin{frame}{Setiap subtes memiliki tekanan waktu yang berbeda}
  \kicker{Jumlah Soal dan Durasi Resmi 2026}
  {\fontsize{9.8}{12.2}\selectfont
  \renewcommand{\arraystretch}{0.92}
  \begin{tabularx}{\textwidth}{>{\RaggedRight\arraybackslash}X >{\centering\arraybackslash}p{0.13\textwidth} >{\centering\arraybackslash}p{0.15\textwidth} >{\centering\arraybackslash}p{0.18\textwidth}}
    \toprule
    \textbf{Subtes} & \textbf{Soal} & \textbf{Menit} & \textbf{Rata-rata/soal} \\
    \midrule
    Penalaran Umum & 30 & 30 & 60 detik \\
    Pengetahuan dan Pemahaman Umum & 20 & 15 & 45 detik \\
    Pemahaman Bacaan dan Menulis & 20 & 25 & 75 detik \\
    Pengetahuan Kuantitatif & 20 & 20 & 60 detik \\
    Literasi dalam Bahasa Indonesia & 30 & 42,5 & 85 detik \\
    Literasi dalam Bahasa Inggris & 20 & 20 & 60 detik \\
    Penalaran Matematika & 20 & 42,5 & 127,5 detik \\
    \bottomrule
  \end{tabularx}}
  \vspace{0.08cm}
  {\color{StudentGray}\fontsize{8.8}{10.8}\selectfont Rata-rata hanya ilustrasi; pembagian waktu tetap disesuaikan dengan kesulitan soal.}
\end{frame}

% Slide 5 ---------------------------------------------------------------------
\begin{frame}{TPS mengukur kesiapan berpikir untuk belajar di perguruan tinggi}
  \kicker{Kelompok 1 -- Tes Potensi Skolastik}
  \begin{columns}[T,onlytextwidth]
    \begin{column}{0.47\textwidth}
      \claim{TPS tidak berfokus pada hafalan mata pelajaran tertentu.}
      \vspace{0.22cm}
      {\fontsize{12}{15}\selectfont Kemampuan yang dilihat berkembang melalui proses belajar dan pengalaman, di dalam maupun di luar kelas.}
    \end{column}
    \begin{column}{0.47\textwidth}
      \begin{enumerate}
        \item Penalaran Umum
        \item Pengetahuan dan Pemahaman Umum
        \item Pemahaman Bacaan dan Menulis
        \item Pengetahuan Kuantitatif
      \end{enumerate}
    \end{column}
  \end{columns}
  \vspace{0.18cm}
  \begin{beamercolorbox}[rounded=true,sep=0.17cm,wd=\textwidth]{block body}
    \centering\bfseries Pembimbingan TPS melatih pola berpikir, pemahaman bahasa, ketelitian, dan keputusan dalam waktu terbatas.
  \end{beamercolorbox}
\end{frame}

% Slide 6 ---------------------------------------------------------------------
\begin{frame}{Penalaran Umum menguji cara menemukan dan memakai aturan}
  \kicker{30 Soal \textbar\ 30 Menit}
  \begin{columns}[T,onlytextwidth]
    \begin{column}{0.43\textwidth}
      {\color{StudentBlue}\bfseries Tiga bagian}\par
      \begin{itemize}
        \item Induktif: menemukan pola dari fakta.
        \item Deduktif: menarik simpulan dari premis.
        \item Kuantitatif: bernalar dengan angka.
      \end{itemize}
      \vspace{0.10cm}
      {\color{StudentGray}\fontsize{10.8}{13.5}\selectfont Masing-masing terdiri atas 10 soal dan 10 menit.}
    \end{column}
    \begin{column}{0.51\textwidth}
      \begin{block}{Contoh singkat: penalaran induktif}
        Perhatikan pola bilangan berikut:
        \[
          5,\;8,\;14,\;23,\;35,\;\ldots
        \]
        Bilangan selanjutnya adalah \ldots
        \answer{50, karena selisihnya bertambah 3: 3, 6, 9, 12, 15.}
      \end{block}
    \end{column}
  \end{columns}
\end{frame}

% Slide 7 ---------------------------------------------------------------------
\begin{frame}{PPU menilai bahasa melalui konteks, bukan hafalan istilah}
  \kicker{20 Soal \textbar\ 15 Menit}
  \begin{columns}[T,onlytextwidth]
    \begin{column}{0.46\textwidth}
      {\color{StudentBlue}\bfseries Materi utama}\par
      \begin{itemize}
        \item Makna kata, frasa, dan kalimat.
        \item Hubungan antargagasan.
        \item Informasi umum dalam konteks sosial budaya.
        \item Simpulan dari bacaan pendek atau panjang.
      \end{itemize}
    \end{column}
    \begin{column}{0.48\textwidth}
      \begin{block}{Contoh singkat}
        "Program belajar harus \textbf{adaptif} terhadap hasil diagnostik siswa."

        Makna kata \textit{adaptif} pada kalimat tersebut adalah \ldots

        \answer{mampu menyesuaikan diri dengan keadaan.}
      \end{block}
    \end{column}
  \end{columns}
\end{frame}

% Slide 8 ---------------------------------------------------------------------
\begin{frame}{PBM menguji kemampuan memperbaiki gagasan tertulis}
  \kicker{20 Soal \textbar\ 25 Menit}
  \begin{columns}[T,onlytextwidth]
    \begin{column}{0.44\textwidth}
      {\color{StudentBlue}\bfseries Tiga cakupan}\par
      \begin{itemize}
        \item Pemilihan gagasan: judul, diksi, pelengkap paragraf.
        \item Penataan gagasan: urutan, posisi, dan penggabungan kalimat.
        \item Perbaikan gagasan: ejaan, keefektifan, dan ketepatan kalimat.
      \end{itemize}
    \end{column}
    \begin{column}{0.50\textwidth}
      \begin{block}{Contoh singkat}
        Kalimat tidak efektif:

        "Karena nilai tryout menurun, sehingga siswa mengikuti remedial."

        \answer{Karena nilai tryout menurun, siswa mengikuti remedial.}
      \end{block}
    \end{column}
  \end{columns}
\end{frame}

% Slide 9 ---------------------------------------------------------------------
\begin{frame}{PK memerlukan konsep matematika dan pilihan strategi yang cepat}
  \kicker{20 Soal \textbar\ 20 Menit}
  \begin{columns}[T,onlytextwidth]
    \begin{column}{0.43\textwidth}
      {\color{StudentBlue}\bfseries Materi utama}\par
      \begin{itemize}
        \item Bilangan dan operasi.
        \item Aljabar.
        \item Geometri dan pengukuran.
        \item Analisis data.
        \item Peluang.
      \end{itemize}
      \vspace{0.08cm}
      {\color{StudentGray}\fontsize{9.8}{12}\selectfont Rumus saja tidak cukup; siswa harus memilih cara yang efisien.}
    \end{column}
    \begin{column}{0.51\textwidth}
      \begin{block}{Contoh singkat}
        Perbandingan peserta kelompok A dan B adalah $3:5$. Jika jumlah seluruh peserta 320 orang, selisih jumlah kedua kelompok adalah \ldots
        Dengan $A=120$ dan $B=200$.\par
        \answer{80 orang.}
      \end{block}
    \end{column}
  \end{columns}
\end{frame}

% Slide 10 --------------------------------------------------------------------
\begin{frame}{Tes Literasi menilai kemampuan menggunakan informasi tertulis}
  \kicker{Kelompok 2 -- Tes Literasi}
  \begin{columns}[T,onlytextwidth]
    \begin{column}{0.45\textwidth}
      \claim{Siswa harus menemukan, memahami, menilai, dan merefleksikan informasi.}
      \vspace{0.20cm}
      {\fontsize{11.8}{14.8}\selectfont Bacaan dan persoalan disajikan dalam berbagai konteks, bukan sebagai pertanyaan hafalan.}
    \end{column}
    \begin{column}{0.49\textwidth}
      \begin{enumerate}
        \item Literasi dalam Bahasa Indonesia
        \item Literasi dalam Bahasa Inggris
        \item Penalaran Matematika
      \end{enumerate}
      \vspace{0.22cm}
      \begin{block}{Kebiasaan yang dibutuhkan}
        Membaca aktif, menandai bukti, membedakan fakta dan simpulan, serta memeriksa kewajaran jawaban.
      \end{block}
    \end{column}
  \end{columns}
\end{frame}

% Slide 11 --------------------------------------------------------------------
\begin{frame}{LBI menguji literasi ilmu pengetahuan dalam tiga konteks}
  \kicker{30 Soal \textbar\ 42,5 Menit}
  \begin{columns}[T,onlytextwidth]
    \begin{column}{0.31\textwidth}
      \begin{block}{Saintek}
        Teks berkaitan dengan fisika, kimia, dan biologi.
      \end{block}
    \end{column}
    \begin{column}{0.31\textwidth}
      \begin{block}{Sosial humaniora}
        Teks berkaitan dengan ekonomi, sejarah, dan sosiogeografi.
      \end{block}
    \end{column}
    \begin{column}{0.31\textwidth}
      \begin{block}{Personal}
        Teks berkaitan dengan apresiasi sastra.
      \end{block}
    \end{column}
  \end{columns}
  \vspace{0.20cm}
  {\color{StudentBlue}\bfseries Kemampuan yang dinilai}\par
  \vspace{0.05cm}
  {\fontsize{11.5}{14.3}\selectfont Menemukan informasi, memahami hubungan gagasan, menilai bukti, menarik simpulan, dan merefleksikan bacaan.}
  \vspace{0.12cm}
  {\color{StudentGray}\fontsize{9.8}{12}\selectfont Domain mengikuti bahan ajar SMA/sederajat; soal tetap berpusat pada literasi.}
\end{frame}

% Slide 12 --------------------------------------------------------------------
\begin{frame}{Skor LBI 2026 membuat profil konteks semakin penting}
  \kicker{Perubahan Hasil dan Antisipasi 2027}
  {\fontsize{10.3}{12.8}\selectfont
  \renewcommand{\arraystretch}{0.98}
  \begin{tabularx}{\textwidth}{>{\RaggedRight\arraybackslash}p{0.24\textwidth} >{\RaggedRight\arraybackslash}p{0.20\textwidth} >{\RaggedRight\arraybackslash}X}
    \toprule
    \textbf{Perubahan} & \textbf{Status} & \textbf{Implikasi pembimbingan} \\
    \midrule
    Sertifikat 2026 memisahkan skor LBI saintek dan soshum & \statusok & Analisis latihan dibedakan menurut konteks bacaan. \\
    Latihan wawasan kebangsaan dalam LBI & \statuswait & Disiapkan sebagai konteks bacaan; menunggu ketentuan SNPMB 2027. \\
    \bottomrule
  \end{tabularx}}
  \vspace{0.12cm}
  \begin{beamercolorbox}[rounded=true,sep=0.14cm,wd=\textwidth]{block body}
    \centering\bfseries Profil konteks menunjukkan jenis bacaan yang masih perlu diperkuat.
  \end{beamercolorbox}
\end{frame}

% Slide 13 --------------------------------------------------------------------
\begin{frame}{Contoh LBI: jawaban harus bertumpu pada bukti bacaan}
  \kicker{Ilustrasi Soal Konteks Saintek}
  \begin{block}{Bacaan singkat}
    Hutan mangrove mengurangi energi gelombang dan menjadi habitat berbagai biota. Ketika mangrove dialihfungsikan, risiko abrasi meningkat dan hasil tangkapan nelayan dapat menurun.
  \end{block}
  \vspace{0.06cm}
  {\fontsize{10.8}{13.5}\selectfont Simpulan yang paling tepat adalah \ldots}
  {\fontsize{10.5}{13}\selectfont
  \begin{enumerate}
    \item Mangrove hanya bermanfaat bagi nelayan.
    \item Alih fungsi mangrove tidak memengaruhi garis pantai.
    \item Pelestarian mangrove mendukung perlindungan pantai dan kehidupan biota.
  \end{enumerate}}
  {\color{StudentGreen}\bfseries\fontsize{10.5}{13}\selectfont Jawaban: C, karena didukung oleh dua fungsi dalam bacaan.}
\end{frame}

% Slide 14 --------------------------------------------------------------------
\begin{frame}{LBE menguji pemahaman bacaan akademik berbahasa Inggris}
  \kicker{20 Soal \textbar\ 20 Menit}
  \begin{columns}[T,onlytextwidth]
    \begin{column}{0.42\textwidth}
      {\color{StudentBlue}\bfseries Konteks bacaan}\par
      \begin{itemize}
        \item Science and technology
        \item Social sciences and humanities
        \item Personal and literary texts
      \end{itemize}
      \vspace{0.10cm}
      {\color{StudentGray}\fontsize{10.5}{13}\selectfont Fokusnya bukan sekadar tata bahasa, tetapi menemukan dan mengevaluasi informasi.}
    \end{column}
    \begin{column}{0.52\textwidth}
      \begin{block}{Contoh singkat}
        "Regular retrieval practice strengthens long-term memory more effectively than repeatedly rereading the same notes."

        What is the main idea?

        \answer{Active recall is more effective for long-term memory than rereading.}
      \end{block}
    \end{column}
  \end{columns}
\end{frame}

% Slide 15 --------------------------------------------------------------------
\begin{frame}{Penalaran Matematika menghubungkan masalah nyata dan model}
  \kicker{20 Soal \textbar\ 42,5 Menit}
  \begin{columns}[T,onlytextwidth]
    \begin{column}{0.47\textwidth}
      {\color{StudentBlue}\bfseries Tiga proses utama}\par
      \begin{enumerate}
        \item \textbf{Merumuskan}: mengubah situasi menjadi model matematika.
        \item \textbf{Menggunakan}: menerapkan konsep dan prosedur.
        \item \textbf{Menafsirkan}: mengembalikan hasil ke konteks masalah.
      \end{enumerate}
    \end{column}
    \begin{column}{0.47\textwidth}
      {\color{StudentBlue}\bfseries Konteks masalah}\par
      \begin{itemize}
        \item Pribadi
        \item Pekerjaan
        \item Sosial
        \item Ilmiah
      \end{itemize}
      \begin{block}{Yang dicari}
        Strategi yang tepat, argumen logis, pola, serta kewajaran hasil.
      \end{block}
    \end{column}
  \end{columns}
\end{frame}

% Slide 16 --------------------------------------------------------------------
\begin{frame}{Contoh Penalaran Matematika: hasil harus masuk akal}
  \kicker{Ilustrasi Pemodelan}
  \begin{block}{Masalah}
    Biaya taksi terdiri atas tarif awal Rp12.000 dan Rp4.000 untuk setiap kilometer. Jika seorang siswa memiliki anggaran paling banyak Rp40.000, berapa kilometer maksimum yang dapat ditempuh?
  \end{block}
  \vspace{0.15cm}
  \begin{columns}[T,onlytextwidth]
    \begin{column}{0.46\textwidth}
      {\color{StudentBlue}\bfseries Model}\par
      \[
        12.000 + 4.000x \leq 40.000
      \]
    \end{column}
    \begin{column}{0.46\textwidth}
      {\color{StudentBlue}\bfseries Interpretasi}\par
      \[
        x \leq 7
      \]
      \answer{Maksimum 7 kilometer.}
    \end{column}
  \end{columns}
\end{frame}

% Slide 17 --------------------------------------------------------------------
\begin{frame}{Profil per subtes lebih berguna daripada satu nilai rata-rata}
  \kicker{Cara Membaca Hasil Tryout}
  \begin{columns}[T,onlytextwidth]
    \begin{column}{0.46\textwidth}
      \begin{block}{Nilai total menjawab}
        Seberapa jauh hasil siswa pada satu simulasi?
      \end{block}
      \vspace{0.15cm}
      \begin{block}{Profil subtes menjawab}
        Mengapa hasil itu terjadi dan bagian mana yang perlu diperbaiki?
      \end{block}
    \end{column}
    \begin{column}{0.48\textwidth}
      \begin{itemize}
        \item Akurasi per subtes dan materi.
        \item Waktu pengerjaan.
        \item Pola jawaban kosong atau tebakan.
        \item Jenis kesalahan yang berulang.
        \item Perubahan hasil dari waktu ke waktu.
      \end{itemize}
    \end{column}
  \end{columns}
\end{frame}

% Slide 18 --------------------------------------------------------------------
\begin{frame}{StudentPro menerjemahkan framework menjadi kegiatan belajar}
  \kicker{Implementasi Menuju 2027}
  {\fontsize{9.5}{11.8}\selectfont
  \renewcommand{\arraystretch}{0.90}
  \begin{tabularx}{\textwidth}{>{\bfseries\color{StudentBlue}}p{0.20\textwidth} >{\RaggedRight\arraybackslash}X >{\RaggedRight\arraybackslash}p{0.29\textwidth}}
    \toprule
    \textbf{Fungsi LMS} & \textbf{Implementasi} & \textbf{Keputusan pembimbingan} \\
    \midrule
    Materi & Modul per subtes dan indikator & Menentukan urutan belajar \\
    Latihan & Kuis topikal dan paket subtes & Mengukur penguasaan bertahap \\
    Tryout & Durasi resmi, urutan, dan simulasi tanpa jeda & Melatih strategi dan stamina \\
    Analisis & Akurasi, waktu, indikator, dan konteks LBI & Menentukan remedial/pengayaan \\
    Monitoring & Ketuntasan, tren, dan konsultasi & Komunikasi sekolah-orang tua \\
    \bottomrule
  \end{tabularx}}
\end{frame}

% Slide 19 --------------------------------------------------------------------
\begin{frame}{Orang tua tidak perlu menguasai semua soal untuk bisa membantu}
  \kicker{Hal yang Perlu Diingat}
  \begin{itemize}
    \item Tanyakan subtes yang sedang diperbaiki, bukan hanya nilai total.
    \item Perhatikan konsistensi belajar dan kondisi fisik saat tryout penuh.
    \item Beri ruang untuk remedial; penurunan sementara dapat menjadi bagian dari proses belajar.
    \item Gunakan laporan StudentPro sebagai bahan dialog, bukan alat membandingkan anak.
    \item Tunggu informasi resmi sebelum mengambil keputusan berdasarkan rumor perubahan UTBK 2027.
  \end{itemize}
  \vspace{0.24cm}
  \begin{beamercolorbox}[rounded=true,sep=0.20cm,wd=\textwidth]{block body}
    \centering\bfseries Bab berikutnya membahas evaluasi pelaksanaan pembimbingan UTBK 2026 di SMAN 1 Muntilan.
  \end{beamercolorbox}
\end{frame}

% Slide 20 --------------------------------------------------------------------
\begin{frame}{Sumber resmi}
  \kicker{Rujukan Bab 2 -- Bagian 1}
  {\color{StudentBlue}\bfseries\fontsize{15}{18}\selectfont 1. Panitia SNPMB}\par
  \vspace{0.06cm}
  {\fontsize{12.5}{15.5}\selectfont "Framework UTBK-SNBT Tahun 2026."}\par
  {\fontsize{11.5}{14}\selectfont\href{https://snpmb.id/fr/}{https://snpmb.id/fr/}}\par
  {\color{StudentGray}\fontsize{11}{13.5}\selectfont Sumber struktur, tujuan, materi, jumlah soal, dan durasi tujuh subtes.}\par
  \vspace{0.38cm}
  {\color{StudentBlue}\bfseries\fontsize{15}{18}\selectfont 2. Panitia SNPMB}\par
  \vspace{0.06cm}
  {\fontsize{12.5}{15.5}\selectfont "Informasi Umum UTBK-SNBT 2026."}\par
  {\fontsize{11.2}{13.8}\selectfont\href{https://snpmb.id/utbk-snbt/informasi-umum}{snpmb.id/utbk-snbt/informasi-umum}}\par
  {\color{StudentGray}\fontsize{11}{13.5}\selectfont Sumber tujuan, ketentuan umum, materi tes, dan pilihan program studi.}\par
  \vspace{0.30cm}
  {\color{StudentGray}\fontsize{10.5}{13}\selectfont Perhitungan total 155 soal dan rata-rata waktu per soal diturunkan dari data resmi per subtes.}
\end{frame}

% Slide 21 --------------------------------------------------------------------
\begin{frame}{Sumber publik dan catatan status}
  \kicker{Rujukan Bab 2 -- Bagian 2}
  {\color{StudentBlue}\bfseries\fontsize{15}{18}\selectfont 3. DetikEdu -- 26 Mei 2026}\par
  \vspace{0.06cm}
  {\fontsize{11}{13.5}\selectfont "Skor LBI di Sertifikat UTBK 2026 Dipecah Saintek-Soshum."}\par
  {\fontsize{10.2}{12.5}\selectfont\href{https://www.detik.com/edu/seleksi-masuk-pt/d-8505258/skor-lbi-di-sertifikat-utbk-2026-dipecah-saintek-soshum-snpmb-jelaskan-mengapa}{detik.com/edu/seleksi-masuk-pt/d-8505258/}}\par
  {\color{StudentGray}\fontsize{10}{12.5}\selectfont Laporan konferensi pers yang mengutip pimpinan SNPMB 2026.}\par
  \vspace{0.24cm}
  {\color{StudentBlue}\bfseries\fontsize{15}{18}\selectfont 4. Rencana integrasi TWK ke LBI}\par
  \vspace{0.06cm}
  {\fontsize{11}{13.5}\selectfont Arah antisipasi materi StudentPro, bukan ketentuan resmi UTBK-SNBT 2027.}\par
  {\color{StudentGray}\fontsize{10}{12.5}\selectfont Status: menunggu framework resmi SNPMB 2027.}\par
  \vfill
  \begin{beamercolorbox}[rounded=true,sep=0.14cm,wd=\textwidth]{block body}
    \centering\bfseries Contoh soal Bab 2 disusun sendiri untuk menjelaskan kemampuan yang dilatih.
  \end{beamercolorbox}
\end{frame}

\end{document}
